% ============================================================
%  Chapter 1 -- Introduction, 1.1-1.4
% ============================================================
\rchapter{1}{The Real and Complex Number Systems}

\begin{roadmap}[Why this chapter exists]
Consider the sequence
\[
  1,\quad 1.4,\quad 1.41,\quad 1.414,\quad 1.4142,\quad \dots
\]
Every number in this sequence is rational. The sequence approaches $\sqrt2$,
but $\sqrt2$ is not rational. Therefore the sequence has no limit that belongs
to $\Q$.

\medskip
\begin{center}
\begin{tikzpicture}[
  every node/.style={font=\scriptsize\sffamily},
  value/.style={draw=RuInk, line width=0.5pt, minimum height=5mm,
                inner xsep=3.5pt, inner ysep=1pt}
]
  % the rational approximations
  \node[value] (a) {$1.4$};
  \node[value, right=2.6mm of a] (b) {$1.41$};
  \node[value, right=2.6mm of b] (c) {$1.414$};
  \node[value, right=2.6mm of c] (d) {$1.4142$};
  \node[right=2mm of d] (dots) {$\cdots$};

  % the place where the limit should be, and is not
  \node[circle, draw=RuInk, dashed, line width=0.5pt, fill=white,
        minimum size=7pt, inner sep=0pt, right=3.5mm of dots] (gap) {};

  \draw[rmarrow] (a) -- (b);
  \draw[rmarrow] (b) -- (c);
  \draw[rmarrow] (c) -- (d);
  \draw[rmarrow] (d) -- (dots);
  \draw[rmarrow] (dots) -- (gap);

  \node[above=2mm of c.north, anchor=south, align=center]
    {every term lies in $\Q$};
  \node[right=1.5mm of gap, align=left]
    {$\sqrt2 \notin \Q$:\\\textbf{limit missing}};

  % the same point, once Q is enlarged
  \node[circle, fill=RuInk, minimum size=7pt, inner sep=0pt,
        below=13mm of gap] (filled) {};
  \draw[-{Stealth[length=5pt]}, line width=0.5pt] (gap) -- (filled);
  \node[left=1.5mm of filled, align=right]
    {enlarge $\Q$\\to $\R$};
  \node[right=1.5mm of filled, align=left]
    {$\sqrt2 \in \R$:\\\textbf{limit present}};
\end{tikzpicture}
\end{center}

\medskip
The real numbers fill this gap by including $\sqrt2$. Rudin secures the property
that rules out such gaps before using it throughout analysis --- he states it as
the \emph{least-upper-bound property} (1.10), and it is what the word
\emph{complete} will name in 3.12.

\smallskip
Everything asserting that something \emph{exists} depends on it: that a bounded
monotone sequence converges (3.14), that a continuous function on $[a,b]$
attains its bounds (4.16) and takes every intermediate value (4.23), that a
continuous function is integrable (6.8). Each fails, or loses its meaning, over $\Q$ --- $x^2-2$
changes sign on $[1,2]\cap\Q$ without ever vanishing. Purely algebraic results
survive; the existence theorems do not.
\end{roadmap}

\begin{roadmap}[How it is organized, and why in that order]
\textbf{Rudin's own account.} The obvious order would be to construct $\R$ from
$\Q$ first and then do analysis on it. He declines, and says why in the Preface:
beginning with the construction is ``pedagogically unsound (though logically
correct),'' because ``most students simply fail to appreciate the need for doing
this.'' Accordingly $\R$ ``is introduced as an ordered field with the
least-upper-bound property,'' and Dedekind's construction is moved to the
appendix, ``where it may be studied and enjoyed whenever the time seems ripe.''

\smallskip
\textbf{Reading that off the chapter}, the shape is one designed to make you
feel the need before paying the cost:

\medskip
\centering
\begin{tikzpicture}[scale=0.92, transform shape, node distance=5mm]
\node[rmbox, text width=22mm] (a) {\textbf{1.1--1.2}\\exhibit the gap\\in $\Q$};
\node[rmbox, text width=22mm, right=of a] (b) {\textbf{1.5--1.18}\\vocabulary:\\order, field};
\node[rmbox, text width=22mm, right=of b] (c) {\textbf{1.19}\\assert $\R$\\(proof deferred)};
\node[rmbox, text width=22mm, right=of c] (d) {\textbf{1.20--1.21}\\cash it out\\at once};
\draw[rmarrow] (a) -- (b); \draw[rmarrow] (b) -- (c); \draw[rmarrow] (c) -- (d);
\end{tikzpicture}

\medskip
\raggedright
Four consequences are worth knowing in advance.

\smallskip
\textbf{The construction is missing on purpose.} Theorem 1.19 is stated and not
proved; nothing in Chapters 2--11 depends on having read the appendix.

\smallskip
\textbf{The abstraction comes before the object.} Ordered sets (1.5--1.11) and
fields (1.12--1.18) are defined \emph{in general} before $\R$ is named, which
lets Rudin characterise $\R$ by properties rather than by construction. He gives
the second half of the reason himself at 1.13(c): proving the elementary
consequences once, for an arbitrary field, means ``we will not need to do it
again for the real numbers and for the complex numbers.''

\smallskip
\textbf{Two of the sections are equipment.} The complex field (1.24--1.35) and
Euclidean space (1.36--1.38) prove nothing about $\R$; they supply objects
Chapters 2--11 need. The one result among them with real downstream weight is
the Schwarz inequality, 1.35, which yields the metric on $\R^k$ and reappears
for integrals in Exercise 6.10 (H\"older's inequality, whose case $p=q=2$ is
there named Schwarz) and for square-integrable functions in 11.35.

\smallskip
\textbf{Where the difficulty lies} --- a reader's judgement, not Rudin's. Not in
1.19, which is asserted, nor in the field axioms, which are routine. It is in
1.21 and in the appendix, the two places where a real number is produced out of
nothing but a supremum. Follow those two arguments and you have the chapter.
\end{roadmap}

\begin{roadmap}[How the proofs in this chapter are found]
A proof you can only verify is a proof you cannot reproduce. Every proof in
this chapter is built from one of seven discovery patterns, catalogued here
once; the box at each proof works the details, and the same patterns carry
most of Chapters 2 and 3.

\smallskip
\textbf{1. Unwind the definition.} When the objects are abstract, the axioms
are the entire arsenal. Ask which axiom can \emph{produce or remove} the
symbol standing between you and the goal, and the proof assembles itself
(1.14, 1.15, 1.16, 1.18).

\smallskip
\textbf{2. Hunt an invariant.} To prove something impossible, find a quantity
the two sides of the supposed equation cannot share --- parity in 1.1, an
integer strictly between $0$ and $1$ in 3.32, strict decrease in Tennenbaum's
dissection.

\smallskip
\textbf{3. Transport the problem to a supremum.} Completeness is the
chapter's only existence tool, so a proof that must \emph{produce} a number
has one move: build a set whose supremum is the number wanted --- $L$ in
1.11, $\{nx\}$ in 1.20(a), $E$ in 1.21 --- then interrogate the supremum with
the two clauses of 1.8: step back to find members, step past to find bounds.

\smallskip
\textbf{4. Work backwards for the constants.} Every mysterious choice ($h$ in
1.21, the coefficients $B$, $C$ in 1.35) is found by writing down the target
inequality and solving for the hypothesis, then presenting the steps in
reverse.

\smallskip
\textbf{5. Prove identities through uniqueness.} To show two expressions are
equal, show both answer a question that has exactly one answer (Corollary to
1.21, 1.33(c)).

\smallskip
\textbf{6. Square away the absolute values.} Statements about $|\cdot|$
become inner-product algebra after squaring; expand, bound the cross term,
take roots at the end (1.33(e), 1.35, 1.37(e)).

\smallskip
\textbf{7. Recognise pure verification.} Some proofs (1.25--1.29, 1.31, 1.37
(a)--(c), most of the appendix) contain no idea to find --- the idea lives in
the \emph{definition} being verified, and the proof is bookkeeping. Knowing
that nothing is hidden is itself the skill: do not search a verification for
insight that was spent earlier, at the definition.
\end{roadmap}

\rsection{Introduction}

\noindent
A satisfactory discussion of the main concepts of analysis (such as convergence,
continuity, differentiation, and integration) must be based on an accurately
defined number concept. We shall not, however, enter into any discussion of the
axioms that govern the arithmetic of the integers, but assume familiarity with
the rational numbers (i.e., the numbers of the form $m/n$ where $m$ and $n$ are
integers and $n \ne 0$).\kn{Rudin is choosing a starting line, and he defends
the choice in the Preface: beginning with a construction of $\R$ is
``pedagogically unsound (though logically correct)'' because ``most students
simply fail to appreciate the need for doing this.'' So the construction is
exiled to the appendix, and $\R$ will arrive in 1.19 as an assertion. This
sentence marks everything he is willing to take for granted.}

The rational number system is inadequate for many purposes, both as a field and
as an ordered set. (These terms will be defined in Definitions 1.6 and 1.12.)
For instance, there is no rational $p$ such that $p^2 = 2$. (We shall prove this
presently.) This leads to the introduction of so-called ``irrational numbers''
which are often written as infinite decimal expansions and are considered to be
``approximated'' by the corresponding finite decimals. Thus the sequence
\[ 1,\ 1.4,\ 1.41,\ 1.414,\ 1.4142,\ \dots \]
``tends to $\sqrt 2$.''\kw{The scare quotes are doing real work. ``Tends to'' is
undefined until Definition 3.1, and the thing it tends to is undefined until
Theorem 1.19. Rudin is pointing at a circularity he is about to break, not
relying on one.} But unless the irrational number $\sqrt 2$ has been clearly
defined, the question must arise: Just what is it that this sequence ``tends
to''?

This sort of question can be answered as soon as the so-called ``real number
system'' is constructed.

\ritem{1.1}{Example}
We now show that the equation
\begin{equation}
  p^2 = 2 \tag{1.1}
\end{equation}
is not satisfied by any rational $p$. If there were such a $p$, we could write
$p = m/n$ where $m$ and $n$ are integers that are not both even.\kn{Why we may
assume this: put $p$ in lowest terms. If $m$ and $n$ were both even we could
cancel a factor of $2$ and repeat, and the process must stop because $|m|$
strictly decreases. That appeal to ``must stop'' is well-ordering, so the step
is not quite free. It can be avoided: Tao (\emph{Analysis I}, 4.4.4) drops the
assumption and instead turns each solution $p^2 = 2q^2$ into a smaller one,
contradicting the principle of infinite descent. Same content, with the
well-ordering moved into the open.} Let us assume this is done. Then (1.1) implies
\begin{equation}
  m^2 = 2n^2, \tag{1.2}
\end{equation}
This shows that $m^2$ is even. Hence $m$ is even (if $m$ were odd, $m^2$ would
be odd), and so $m^2$ is divisible by $4$. It follows that the right side of
(1.2) is divisible by $4$, so that $n^2$ is even, which implies that $n$ is
even.

The assumption that (1.1) holds thus leads to the conclusion that both $m$ and
$n$ are even, contrary to our choice of $m$ and $n$. Hence (1.1) is impossible
for rational $p$.

\begin{unpack}[How the parity argument is found]
Faced with $m^2 = 2n^2$, the search is for an \emph{invariant}: some property
that the equation forces on both sides but that the two sides cannot share.
The right side is visibly even, so ask what evenness of $m^2$ says about $m$
--- and the cheapest arithmetic invariant, parity, turns out to propagate:
$m$ even makes $m^2$ divisible by $4$, pushing evenness onto $n^2$ and then
$n$.

Two design choices become visible once you look for them. The normalisation
``not both even'' exists so that the invariant has something to contradict ---
without it, ``both even'' is not absurd, merely reducible. And the argument is
really about the exponent of $2$: each side of $m^2 = 2n^2$ contains the prime
$2$ an even number of times on the left and an odd number on the right, which
is the version that generalises (Exercise 1.2 runs it with the prime $3$ on
$m^2 = 12n^2$). Parity is that count, read modulo $2$.
\end{unpack}

\begin{deeper}[An alternative proof, done entirely by cutting up a square]
Rudin argues by parity. There is a second proof, due to Stanley Tennenbaum,
which uses no arithmetic beyond areas --- and it makes the descent visible
rather than hidden in the phrase ``not both even.''

\medskip
Suppose $m^2 = 2n^2$ with $m,n$ positive integers. Then $n<m<2n$, since
$m^2=2n^2$ lies strictly between $n^2$ and $4n^2$. Draw a square of side $m$
and lay a square of side $n$ into two opposite corners. Because $m<2n$ they
overlap; because $m>n$ two corners are left bare.

\medskip
\begin{center}
\begin{tikzpicture}[x=1mm, y=1mm]
  % the two uncovered corner squares
  \fill[pattern=north east lines, pattern color=RuInk] (0,24) rectangle (10,34);
  \fill[pattern=north east lines, pattern color=RuInk] (24,0) rectangle (34,10);
  % the doubly covered overlap
  \fill[black!30] (10,10) rectangle (24,24);
  % the two n-squares
  \draw[thin] (0,0)   rectangle (24,24);
  \draw[thin] (10,10) rectangle (34,34);
  % the big square
  \draw[draw=RuInk, line width=0.8pt] (0,0) rectangle (34,34);
  % labels
  \node[lbl, anchor=north] at (17,-2) {side $m$};
  \node[lbl, anchor=east, rotate=90] at (-2,12) {side $n$};
  \node[mlbl] at (17,17) {$2n\!-\!m$};
  \node[mlbl, fill=white, inner sep=1pt] at (5,29)  {$m\!-\!n$};
  \node[mlbl, fill=white, inner sep=1pt] at (29,5)  {$m\!-\!n$};
  % legend
  \fill[black!30] (44,26) rectangle (50,31); \draw[thin] (44,26) rectangle (50,31);
  \node[lbl, anchor=west] at (52,28.5) {covered twice};
  \fill[pattern=north east lines, pattern color=RuInk] (44,17) rectangle (50,22);
  \draw[thin] (44,17) rectangle (50,22);
  \node[lbl, anchor=west] at (52,19.5) {covered not at all};
\end{tikzpicture}
\end{center}
\figcap{Two squares of side $n$ dropped into opposite corners of a square of
side $m$.}

\smallskip
The two $n$-squares have total area $2n^2$, which by hypothesis is exactly
$m^2$, the area of the big square. So whatever is counted twice must have the
same area as whatever is missed:
\[ (2n-m)^2 \;=\; 2\,(m-n)^2 . \]
That is the same equation again, in strictly smaller positive integers ---
$0<2n-m<m$ and $0<m-n<n$. Repeating produces an infinite strictly decreasing
sequence of positive integers, which is impossible. Hence no such $m,n$ exist.

\smallskip
The two proofs use the same principle. Rudin's ``not both even'' silently
invokes well-ordering to stop a descent; here the descent is performed in the
open, and each step is a picture.
\end{deeper}

We now examine this situation a little more closely. Let $A$ be the set of all
positive rationals $p$ such that $p^2 < 2$ and let $B$ consist of all positive
rationals $p$ such that $p^2 > 2$. We shall show that \emph{$A$ contains no
largest number and $B$ contains no smallest}.\kn{This, and not the irrationality
of $\sqrt2$, is the real content of \S1.1. Irrationality says
$\sqrt2 \notin \Q$. This says $\Q$ has a \emph{gap} where $\sqrt 2$ should be
--- a structural defect, and the one Definition 1.10 will name.}

More explicitly, for every $p$ in $A$ we can find a rational $q$ in $A$ such
that $p < q$, and for every $p$ in $B$ we can find a rational $q$ in $B$ such
that $q < p$.

To do this, we associate with each rational $p > 0$ the number
\begin{equation}
  q = p - \frac{p^2-2}{p+2} = \frac{2p+2}{p+2}. \tag{1.3}
\end{equation}
Then
\begin{equation}
  q^2 - 2 = \frac{2(p^2-2)}{(p+2)^2}. \tag{1.4}
\end{equation}
If $p$ is in $A$, then $p^2 - 2 < 0$, (1.3) shows that $q > p$ and (1.4) shows
that $q^2 < 2$. Thus $q$ is in $A$.

If $p$ is in $B$, then $p^2 - 2 > 0$, (1.3) shows that $0 < q < p$, and (1.4)
shows that $q^2 > 2$. Thus $q$ is in $B$.

\begin{unpack}[Where (1.3) comes from]
Rudin produces this formula with no derivation, and it is the first place in the
book where a reader stalls. It is not Newton's method. Ask instead: which map
$p \mapsto q$ has $\sqrt2$ as a \emph{fixed point}? Try the simplest shape, a
ratio of linear terms $q = (ap+b)/(cp+d)$. Setting $q = p$ gives
$p(cp+d) = ap+b$, and we want that to collapse to $p^2 = 2$. The
smallest-integer solution is $q = (2p+2)/(p+2)$: it gives $p^2 + 2p = 2p+2$,
that is, $p^2 = 2$.

Why it moves $p$ toward $\sqrt2$ without crossing it becomes visible once you
factor the error:
\[
  q - \sqrt2 \;=\; \frac{2-\sqrt2}{p+2}\,\bigl(p - \sqrt2\bigr).
\]
For $p>0$ the multiplier lies strictly between $0$ and
$(2-\sqrt2)/2 \approx 0.29$. That it is \emph{positive} keeps $q$ on the same
side of $\sqrt2$ as $p$ --- Rudin's ``$q$ is in $A$'' and ``$q$ is in $B$.'' That
it is \emph{less than one} makes the error shrink, so $q$ is strictly better
than $p$ and no member of $A$ can be largest.

Equation (1.4) is worth checking rather than accepting:
\[ (2p+2)^2 - 2(p+2)^2 = (4p^2+8p+4)-(2p^2+8p+8) = 2p^2-4. \]

\smallskip
There is a second way to say what this establishes, worth knowing because it is
the form other books use. Tao (\emph{Analysis I}, 4.4.5) proves: for every
rational $\varepsilon>0$ there is a rational $x \ge 0$ with
$x^2 < 2 < (x+\varepsilon)^2$. Rudin's version is about the sets $A$ and $B$
having no extreme member; Tao's is about how narrowly they can be straddled.
They say the same thing --- the gap is real, and it is narrower than any
rational you care to name.
\end{unpack}

\begin{center}
\begin{tikzpicture}[x=1mm, y=1mm]
  \draw[axis] (-1,0) -- (94,0);
  \node[ptfill, minimum size=2.6pt] at (4,0) {};
  \node[mlbl, anchor=north] at (4,-1.6) {$0$};
  % A
  \draw[seg] (4,0) -- (44,0);
  \node[lbl, anchor=north] at (24,-2) {$A$: $p^2<2$, no largest};
  % gap
  \node[ptgap] at (47,0) {};
  \node[lbl, anchor=south] at (47,10) {$\sqrt2 \notin \Q$};
  \draw[guide] (47,1.5) -- (47,9.5);
  % B
  \draw[seg] (50,0) -- (88,0);
  \node[lbl, anchor=north] at (69,-2) {$B$: $p^2>2$, no smallest};
  % the map pushes p rightward inside A
  \node[ptfill, minimum size=2.6pt] at (24,0) {};
  \node[ptfill, minimum size=2.6pt] at (34,0) {};
  \node[mlbl, anchor=south] at (24,1.6) {$p$};
  \node[mlbl, anchor=south] at (34,1.6) {$q$};
  \draw[thin, -{Stealth[length=3.5pt]}] (25.5,4) .. controls (28,7) .. (33,4);
  \node[lbl, anchor=south] at (29,7) {$q=\dfrac{2p+2}{p+2}$};
\end{tikzpicture}
\end{center}
\figcap{Every $p\in A$ is beaten by some $q\in A$, and every $p\in B$ beats some
$q\in B$. The two sets crowd a point that is not there.}

\begin{deeper}[The estimate behind the contraction principle]
That factorisation says $p \mapsto (2p+2)/(p+2)$ shrinks the distance to
$\sqrt2$ by a factor of at most $0.29$ at every step. It is the same estimate
that Theorem 9.23, the contraction principle, later abstracts --- though 9.23
itself does not apply here, since it needs a \emph{complete} space and the
positive rationals are exactly what is not complete. That is the point: the
limit $\sqrt2$ does not exist in $\Q$, which is why Rudin must argue by hand
that no member of $A$ is largest rather than simply passing to a fixed point.

Iterating from $p_0 = 1$ gives
\[ 1,\ \tfrac43,\ \tfrac75,\ \tfrac{24}{17},\ \tfrac{41}{29},\
   \tfrac{140}{99},\ \dots \]
Every term stays below $\sqrt2$, as the positive multiplier requires. Some are
convergents of $\sqrt 2 = [1;2,2,2,\dots]$ --- $\tfrac75$ and $\tfrac{41}{29}$
are --- but not all: the convergents run
$1, \tfrac32, \tfrac75, \tfrac{17}{12}, \tfrac{41}{29}, \dots$, and this
iteration visits every second one.

Newton's method on $f(p) = p^2-2$ would give $p \mapsto (p^2+2)/(2p)$ instead:
faster, but from any starting point it lands \emph{above} $\sqrt2$, so it would
prove ``$B$ has no smallest'' and say nothing about $A$. Rudin's map does both
halves at once.
\end{deeper}

\ritem{1.2}{Remark}
The purpose of the above discussion has been to show that the rational number
system has certain gaps, in spite of the fact that between any two rationals
there is another: if $r < s$, then $r < (r+s)/2 < s$.\kn{The property Rudin
describes here --- between any two elements lies a third --- is not the property
Definition 1.10 will demand, and this sentence is the first warning that the two
are independent. $\Q$ has the first and lacks the second. $\Z$ has the second
(every nonempty bounded-above set of integers has a largest element) and lacks
the first. So no amount of ``no visible gaps between neighbours'' delivers
completeness.} The
real number system fills these gaps. This is the principal reason for the
fundamental role which it plays in analysis.

In order to elucidate its structure, as well as that of the complex numbers, we
start with a brief discussion of the general concepts of \emph{ordered set} and
\emph{field}.

Here is some of the standard set-theoretic terminology that will be used
throughout this book.

\ritem{1.3}{Definitions}
If $A$ is any set (whose elements may be numbers or any other objects), we write
$x \in A$ to indicate that $x$ is a member (or an element) of $A$. If $x$ is not
a member of $A$, we write $x \notin A$.

The set which contains no element will be called the \emph{empty set}. If a set
has at least one element, it is called \emph{nonempty}.

If $A$ and $B$ are sets, and if every element of $A$ is an element of $B$, we
say that $A$ is a \emph{subset} of $B$, and we write $A \subset B$, or
$B \supset A$.\kw{Rudin's $\subset$ permits equality --- it is what most modern
texts write $\subseteq$. He writes ``proper subset'' when he means strict
containment. Read every $\subset$ in this book as $\subseteq$, or you will
misread hypotheses for the next 300 pages.} If, in addition, there is an element
of $B$ which is not in $A$, then $A$ is said to be a \emph{proper subset} of
$B$. Note that $A \subset A$ for every set $A$.

If $A \subset B$ and $B \subset A$, we write $A = B$. Otherwise, $A \ne B$.

\ritem{1.4}{Definition}
Throughout Chap.\ 1, the set of all rational numbers will be denoted by $\Q$.
